\section{Sửa bằng cấu trúc nhân quả}
\label{sec:ch14-sua-bang-nhan-qua}

\index{Causal SHAP|(}%
Cả hai phép tấn công vừa đọc đều đi qua cùng một bước: để có điểm số SHAP thì
phải hỏi mô hình về những đầu vào có đặc trưng bị thay hoặc bị che, và những
đầu vào ấy không phải dữ liệu thật. Ở mục~\ref{sec:ch14-phep-dung-gian} đó là
cái khe hở bị khai thác; ở mục~\ref{sec:ch14-shap-lam-vu-khi} đó chỉ là chi phí
phải trả để tính điểm số. Bài báo cuối của chương nhắm vào chính bước ấy, và
nhắm vào nó mà không có bên tấn công nào trong câu chuyện.

\index{Causal SHAP!chỗ hỏng nó sửa}%
Chỗ hỏng mà bài nêu bắt đầu từ một giả định mà chương~\ref{ch:shap} đã ghi và
chưa gỡ. Muốn tính giá trị của một liên minh thiếu vài đặc trưng thì phải điền
giá trị vào những chỗ thiếu, và cách điền phổ biến nhất là lấy chúng độc lập
với những đặc trưng còn lại. Khi các đặc trưng thật sự phụ thuộc nhau, cách
điền ấy dựng ra những điểm dữ liệu không thể tồn tại, và bài lấy một ví dụ y
khoa để nói điều đó cụ thể: hút thuốc gây ung thư phổi, hút thuốc cũng khiến
người ta uống nhiều cà phê hơn, nên uống cà phê tương quan với ung thư phổi mà
không gây ra nó. SHAP thường chia điểm cho cả hút thuốc lẫn uống cà
phê~\autocite{p29causalshap}.

\index{Causal SHAP!hai bước sửa}%
Phép sửa gồm hai bước cộng một bước dọn dẹp, và cả ba đều cần một đồ thị nhân
quả mà bài không giả định là có sẵn mà đi tìm từ dữ liệu. Việc tìm ấy do thuật
toán Peter-Clark, viết tắt PC, đảm nhận: nó bắt đầu từ đồ thị đầy đủ rồi lần
lượt bỏ cạnh giữa những cặp biến độc lập có điều kiện, sau đó định hướng phần
cạnh mà dữ liệu định hướng được~\autocite{p29causalshap}. Sách viết đồ thị ấy
là $G$ và gọi các đỉnh, các cạnh bằng lời, vì $E$ ở phụ lục~\ref{app:ky-hieu}
đã thuộc về tập đối tượng mà một \tn{hàm attribution}{attribution function}
chấm điểm.

Bước thứ nhất đổi cách điền. Thay vì lấy các đặc trưng ngoài liên minh độc lập,
bài lấy chúng theo phân phối mà cấu trúc nhân quả quy định khi các đặc trưng
trong liên minh đã bị can thiệp. Gọi cha của một biến là những biến có cạnh
hướng thẳng vào nó trong $G$, thì quy tắc là: biến nào không có cha thì giữ
nguyên phân phối của riêng nó, còn biến nào có cha thì được sinh ra theo giá
trị mà các cha của nó đang nhận. Hàm giá trị thu được, bài viết là $v_c(S)$, vì
thế không sinh ra những điểm không thể tồn tại~\autocite{p29causalshap}. Đây là
cùng một $v(S)$ của chương~\ref{ch:shap} với một quy tắc điền khác, chứ không
phải một đại lượng mới, nên sách giữ chữ cũ và thêm chỉ số dưới.

\index{cường độ nhân quả}%
Bước thứ hai chiết khấu điểm số theo
\tn{cường độ nhân quả}{causal strength}. Thuật toán IDA, viết tắt
của intervention calculus when the DAG is absent, ước lượng tổng hiệu ứng nhân
quả $W_i$ của từng đặc trưng lên biến mục tiêu, và bước ấy cần một thứ mà PC
không cho: PC trả về đồ thị chứ không trả về độ mạnh của từng cạnh. IDA lấp chỗ
ấy bằng hồi quy, ước lượng trọng số của mỗi cạnh từ dữ liệu, rồi nhân các trọng
số dọc theo từng đường đi từ đặc trưng tới mục tiêu và cộng các đường lại. Hai
thuật toán vì thế chạy nối nhau chứ không độc lập. Bài đặt một hệ số tỉ lệ với
$\lvert W_i\rvert$ đã chuẩn hóa và nhân nó vào trọng số Shapley, nên một đặc
trưng không có đường đi nào tới mục tiêu nhận hệ số bằng
không~\autocite{p29causalshap}. Sách viết hệ số ấy là $\lambda_i$; bài viết
$\gamma_i$, chữ mà chương~\ref{ch:khung-hop-nhat} đã cấp cho \tn{attribution
thành phần}{component attribution}.

Phép nhân ấy phá mất độ chính xác cục bộ, nên bước dọn dẹp là chuẩn hóa lại
tổng các điểm số về đúng hiệu giữa dự đoán tại $x$ và dự đoán trung
bình~\autocite{p29causalshap}. Bước chuẩn hóa này nhỏ và dễ lướt qua, nhưng nó
là chỗ mà phần đánh giá lý thuyết của bài mắc, nên nó được nêu ra ở đây trước
khi tới đó.

\index{Causal SHAP!ba định lý}%
Bài chứng minh rằng ba tính chất của định nghĩa~\ref{def:ch04-ba-tinh-chat} vẫn
còn. Hai trong ba chứng minh đứng vững và một thì không, nên phải tách ra.

Độ chính xác cục bộ đúng, và đúng vì bước chuẩn hóa đặt nó vào bằng định nghĩa:
chứng minh của bài là một dòng ước lược~\autocite{p29causalshap}. Tính chất ấy
vì thế không phải một tính chất mà phép dựng giành được, mà là một tính chất
được áp đặt, và đọc nó như bằng chứng rằng phép dựng đúng là đọc sai.

\index{Causal SHAP!định lý tính vắng mặt}%
\tn{Tính vắng mặt}{missingness} thì bài phát biểu một đằng và chứng minh một nẻo. Ở phần nêu ba
tính chất, bài viết tính vắng mặt theo nghĩa quen thuộc, rằng một đặc trưng
thiếu giá trị phải nhận điểm số bằng không. Định lý mà bài chứng minh nói rằng
một đặc trưng không nằm trong đồ thị nhân quả, hoặc nằm trong đồ thị mà không
có đường đi tới biến mục tiêu, thì nhận điểm số bằng
không~\autocite{p29causalshap}. Chứng minh ấy đúng với điều nó phát biểu, và
điều nó phát biểu không phải tính vắng mặt: vắng mặt khỏi đồ thị nhân quả không
phải vắng mặt của một giá trị đặc trưng. Sách ghi cả hai và không nhận rằng
tiên đề đã được giữ lại.

\index{Causal SHAP!tính nhất quán không sống sót phép chuẩn hóa}%
\tn{Tính nhất quán}{consistency} là chỗ nặng nhất. Bài chứng minh nó cho điểm số trước chuẩn hóa,
rồi viết một câu rằng dễ thấy hai định lý sau cũng đúng cho điểm số đã chuẩn
hóa~\autocite{p29causalshap}. Với tính vắng mặt thì đúng thật, vì kết luận là
số không và số không sống sót mọi phép nhân với số dương. Với tính nhất quán
thì không. Tính nhất quán là một phát biểu so hai mô hình $f$ và $f'$, còn hệ
số chuẩn hóa lại phụ thuộc vào mô hình, vì nó chứa $f(x)$ và tổng các điểm số
chưa chuẩn hóa của chính mô hình ấy. Hai vế của bất đẳng thức vì thế bị nhân
với hai hệ số khác nhau, nên bất đẳng thức trước chuẩn hóa không kéo theo bất
đẳng thức sau chuẩn hóa. Nói gọn lại: đúng bước mua được độ chính xác cục bộ là
bước làm mất tính nhất quán, và bài không thấy chỗ đánh đổi ấy vì nó gộp hai
định lý vào một câu.

\index{Causal SHAP!ground truth tổng hợp}%
Phần thí nghiệm chia làm hai nửa, và nửa thứ nhất là đường dựng ground truth
thứ hai mà mục~\ref{sec:ch14-ground-truth-ke-tan-cong} đã hẹn. Hai bộ dữ liệu
tổng hợp, một bộ về nguy cơ ung thư phổi và một bộ về nguy cơ tim mạch, đều do
bài tự sinh ra từ các phương trình mà bài viết ra: nguy cơ ung thư phổi bằng
hai lần mức hút thuốc cộng 1.2 lần mức căng thẳng cộng nhiễu, và ở bộ kia thì
chỉ số khối cơ thể đứng giữa như một biến gây nhiễu chung. Quan hệ nhân quả
trong dữ liệu vì thế là quan hệ do chính các tác giả đặt
vào~\autocite{p29causalshap}. Giá trị tham chiếu để so thì lấy bằng cách chạy
SHAP trên một tập đặc trưng đã rút gọn, chỉ gồm những đặc trưng có đường nhân
quả trực tiếp tới mục tiêu~\autocite{p29causalshap}.

Chỗ ấy đáng dừng lại một câu. Cái được gọi là ground truth ở đây không phải một
đại lượng độc lập với phương pháp đang được chấm; nó là kết quả của chính SHAP,
chạy trên tập đặc trưng mà người thí nghiệm đã chọn sẵn là tập đúng. Vậy phép
đánh giá hỏi một câu hẹp hơn câu nó có vẻ hỏi: nó hỏi mỗi biến thể có tiến gần
tới SHAP-trên-những-đặc-trưng-đúng hay không, chứ không hỏi biến thể ấy có phản
ánh đúng mô hình hay không.

\index{Causal SHAP!cột nhãn sai}%
Năm mốc so sánh mà bài chọn đều là biến thể của SHAP chứ không có LIME:
Independent SHAP, KernelSHAP, On Manifold SHAP, Shapley Flow và
ASV~\autocite{p29causalshap}. Trong phạm vi câu hỏi hẹp ở trên thì Causal SHAP
thắng, và cột số dùng để tuyên bố điều đó có một chỗ sai mà sách kiểm ra được.

Bài in một cột nhãn là RMSE, tức căn của trung bình bình phương sai lệch so với
giá trị tham chiếu, và Causal SHAP đạt 0.0167 trên bộ ung thư phổi so với
1.6357 của Independent SHAP đứng ngay sau~\autocite{p29causalshap}. Dựng lại
phép tính từ chính các ô trong bảng thì con số ấy không phải một RMSE mà là
tổng bình phương sai lệch, chưa chia cho số đặc trưng và chưa lấy căn. Trên bộ
ung thư phổi, hai đặc trưng có giá trị tham chiếu là hút thuốc và căng thẳng.
Independent SHAP chấm chúng 3.9467 và 0.1031 so với hai giá trị tham chiếu
5.2171 và 0.2507, nên hai sai lệch là $-1.2704$ và $-0.1476$; tổng bình phương
của chúng bằng 1.6357, đúng tới chữ số thập phân thứ tư của ô đã in, trong khi
căn của trung bình bình phương thì bằng 0.9044 và không khớp ô nào. Cách đọc ấy
dựng lại mười một trong mười hai ô của bảng, mỗi ô đúng trong phạm vi sai số mà
bốn chữ số thập phân cho phép sau khi bình phương, và hai ô thì đúng trọn bốn
chữ số. Mười một lần khớp như vậy không còn là trùng hợp.

Đúng một ô không dựng lại được: Shapley Flow ở khối ung thư phổi, bảng in
8.8037 còn phép tính cho 9.3459. Sách nêu nó ra và dừng ở đó, vì mọi cách đọc
khác đã thử đều không cho ra con số của bảng. Chỗ ấy không kéo theo gì cả: cái
nhãn sai không đổi kết luận nào, bởi tổng bình phương và căn của trung bình
bình phương là hai hàm đồng biến của nhau, nên mọi thứ tự trong bảng giữ nguyên
dưới cả hai cách đọc và Causal SHAP vẫn thấp nhất ở cả hai bộ. Đây đúng là
trường hợp mà chương~\ref{ch:gioi-han-ly-thuyet} đã phải học một lần: một phép
suy luận hỏng cho phép bác phép suy luận ấy chứ không cho phép bác kết luận.

\index{insertion!dùng để chấm điểm phép sửa}%
Nửa thứ hai của phần thí nghiệm chạy trên hai bộ dữ liệu y sinh thật, một bộ về
hội chứng ruột kích thích và một bộ về ung thư đại trực tràng, và ở đó không
còn giá trị tham chiếu nào, nên bài phải mượn một dụng cụ đo. Dụng cụ nó mượn
là insertion: thêm dần các đặc trưng từ quan trọng nhất tới ít quan trọng nhất
theo thứ hạng mà mỗi phương pháp đưa ra, rồi đo AUROC, cross entropy và điểm
Brier ở từng bước~\autocite{p29causalshap}. Causal SHAP đạt AUROC 0.8594 trên
bộ thứ nhất và 0.6271 trên bộ thứ hai, cao nhất ở cả hai. Trên bộ thứ hai nó
cũng đứng đầu ở hai chỉ số còn lại; trên bộ thứ nhất thì On Manifold SHAP hơn
nó một chút ở cả cross entropy lẫn điểm Brier, chênh nhau ở chữ số thập phân
thứ tư~\autocite{p29causalshap}.

Bài gọi insertion là một chỉ số chuẩn dùng để đánh giá attribution đặc
trưng~\autocite{p29causalshap}, và chữ \enquote{chuẩn} ấy đúng theo nghĩa nhiều
người dùng. Nó không đúng theo nghĩa đã được kiểm định, và cả
chương~\ref{ch:do-do-trung-thuc} lẫn
chương~\ref{ch:chi-so-khong-do-do-trung-thuc} tồn tại để tách hai nghĩa ấy ra.
Vậy phép sửa duy nhất trong chương này, ở nửa thí nghiệm duy nhất chạy trên dữ
liệu thật, được chấm điểm bằng đúng loại dụng cụ mà Phần~IV vừa dành sáu chương
để nói là chưa ai kiểm.

\index{Causal SHAP|)}%
